\newif\ifgenesissequencestandalone
\ifnum\pdfstrcmp{\jobname}{genesis_sequence}=0
  \genesissequencestandalonetrue
\else
  \genesissequencestandalonefalse
\fi

\ifgenesissequencestandalone
\documentclass[11pt,a4paper]{article}

\usepackage[utf8]{inputenc}
\usepackage[T1]{fontenc}
\usepackage{amsmath,amsthm,amssymb}
\usepackage{mathtools}
\usepackage{array}
\usepackage{enumitem}
\usepackage{booktabs}
\usepackage{geometry}
\usepackage{hyperref}
\usepackage{cleveref}
\usepackage{float}
\usepackage{microtype}
\usepackage{xcolor}
\usepackage{stmaryrd}

\microtypesetup{expansion=false}
\geometry{margin=1in}
\hypersetup{
    colorlinks=true,
    linkcolor=blue!70!black,
    citecolor=green!50!black,
    urlcolor=blue!70!black,
    hypertexnames=false
}
\emergencystretch=2em
\pdfstringdefDisableCommands{%
  \def\leq{<=}%
  \def\geq{>=}%
  \def\to{->}%
  \def\otimes{otimes}%
  \def\nabla{nabla}%
  \def\wedge{wedge}%
  \def\flat{flat}%
  \def\sharp{sharp}%
  \def\bigcirc{next}%
  \def\Diamond{eventually}%
}

\theoremstyle{plain}
\newtheorem{theorem}{Theorem}[section]
\newtheorem{lemma}[theorem]{Lemma}
\newtheorem{proposition}[theorem]{Proposition}
\newtheorem{corollary}[theorem]{Corollary}

\theoremstyle{definition}
\newtheorem{definition}[theorem]{Definition}
\newtheorem{example}[theorem]{Example}
\newtheorem{axiom}[theorem]{Axiom}

\theoremstyle{remark}
\newtheorem{remark}[theorem]{Remark}
\newtheorem{observation}[theorem]{Empirical Observation}

\newcommand{\N}{\mathbb{N}}
\newcommand{\Z}{\mathbb{Z}}
\newcommand{\R}{\mathbb{R}}
\newcommand{\D}{\mathbb{D}}
\newcommand{\U}{\mathcal{U}}
\newcommand{\Disc}{\mathrm{Disc}}

\begin{document}
\fi

\section{A Detailed Proof of the 15-Step Genesis Sequence}
\label{sec:genesis-sequence-proof}

This section gives a step-by-step proof of the current 15-step PEN genesis
sequence in a form intended for a journal proof spine rather than for a
developer report.
The global mathematical input comes from the established PEN theorems in
\cref{thm:scaling,ax:selection,def:derivability,def:minimal-api,thm:end-of-history},
while the step-local selection facts come from the verified strict trace
generated by
\texttt{cabal run ab-initio -- --strict --max-steps 15 --skip-validation --skip-mcts}.
By \cref{rem:automation-scope}, that strict run is the canonical source for the
selection ordering, the bar clearances, and the quantitative values
$(\nu,\kappa,\rho)$ used below.

% Keep proof headings input-safe across LaTeX setups without changing the host
% document's sectioning semantics after this file ends.
\let\genesissequenceorigparagraph\paragraph
\renewcommand{\paragraph}[1]{\par\medskip\noindent\textbf{#1.}\ }

\subsection{Proof Contract}
\label{subsec:genesis-proof-contract}

Three conventions govern the proof.

\begin{enumerate}[nosep]
    \item \textbf{Charged cost.} Throughout this section, $\kappa$ means the
    charged conceptual/desugared cost used by the strict selector, not the raw
    size of every compiled auxiliary clause.
    This is the cost notion fixed in \cref{def:effort,ax:selection} and exposed
    operationally in the strict engine description of
    \cref{sec:verification}.
    \item \textbf{Compiled payload.} Especially for the HIT steps, the compiled
    telescope may contain framework-derived eliminators or computation payload
    beyond the charged conceptual shell.
    Those extra compiled clauses can increase the evaluated novelty, but they do
    not increase the charged $\kappa$.
    This is the conceptual/compiled separation described in
    \cref{sec:verification}.
    \item \textbf{Uniqueness.} Every uniqueness claim below means
    ``unique winner in the current strict unguided lane under admissibility,
    semantic minimality, and minimal overshoot''.
    It does \emph{not} assert a foundation-independent uniqueness theorem over
    all logics or all search policies.
\end{enumerate}

Two further scope boundaries matter for the late stages.
First, the codebase contains additional family-specific recovery and quality
machinery in non-canonical lanes, but the present section does not use those
heuristics as evidence for the strict claim.
Whenever the verified strict trace already yields a singleton viable set, we use
that stronger fact directly.
Second, the proof-strength boundary from \cref{rem:automation-scope} is
maintained verbatim:
steps~1--10 admit exact or exact-reconstruction language,
steps~11--14 are expository reconstructions constrained to match the strict
totals, and step~15 canonizes only the strict total $\nu=88$.

\subsection{A Global Bar Lemma}
\label{subsec:genesis-bar-lemma}

\begin{lemma}[Bar computation along the strict genesis trace]
\label{lem:genesis-bar-lemma}
For every step $n \ge 2$ in the strict genesis trace,
\[
  \mathrm{Bar}_n
  =
  \Phi_n \Omega_{n-1}
  =
  \frac{F_n}{F_{n-1}}
  \cdot
  \frac{\sum_{i=1}^{n-1}\nu_i}{\sum_{i=1}^{n-1}\kappa_i}.
\]
Consequently the active bars for steps~$2$ through~$15$ are exactly the values
listed in \cref{tab:genesis-sequence-trace}.
\end{lemma}

\begin{proof}
By \cref{thm:scaling}, the current $d=2$ lane satisfies $\Delta_n=F_n$, hence
\[
  \Phi_n = \Delta_n/\Delta_{n-1} = F_n/F_{n-1}.
\]
By \cref{def:omega},
\[
  \Omega_{n-1} = \frac{\sum_{i=1}^{n-1}\nu_i}{\sum_{i=1}^{n-1}\kappa_i}.
\]
By \cref{ax:selection}, the selection bar is
$\mathrm{Bar}_n = \Phi_n \Omega_{n-1}$.
Substituting the already realized strict-trace values
$(\nu_i,\kappa_i)$ for $i<n$ therefore determines each bar uniquely.
The worked arithmetic in \cref{app:worked-arithmetic} checks this explicitly for
steps~5 and~6, and the same substitution yields the full bar column in
\cref{tab:genesis-sequence-trace}.
\end{proof}

\begin{table}[tbp]
\centering
\caption{Strict genesis trace used in the proof.
The `cap/bands' column reports the strict admissibility cap together with the
searched exact $\kappa$-band(s).
The `raw/dist.' column reports raw and deduplicated candidate counts.
`Runner-up' and `near miss' mean the top displayed strict-lane competitor in the
bar-clearing and subcritical pools respectively; a dash means that no distinct
competitor was displayed by the verified strict trace.}
\label{tab:genesis-sequence-trace}
\tiny
\setlength{\tabcolsep}{2pt}
\begin{tabular}{@{}c l r r r r r r l l >{\raggedright\arraybackslash}p{1.6cm} >{\raggedright\arraybackslash}p{1.6cm}@{}}
\toprule
Step & Structure & $\Delta_n$ & $\nu$ & $\kappa$ & $\rho$ & Bar & Margin & cap/bands & raw/dist. & Runner-up & Near miss \\
\midrule
1  & Universe    & 1   & 1  & 2 & 0.50 & 0.50 & 0.00 & 2 / 2   & 2144 / 552 & --- & --- \\
2  & Unit        & 1   & 1  & 1 & 1.00 & 0.50 & 0.50 & 1 / 1   & 96 / 78    & --- & --- \\
3  & Witness     & 2   & 2  & 1 & 2.00 & 1.33 & 0.67 & 1 / 1   & 96 / 81    & --- & found.\ cand.\ (1.00) \\
4  & $\Pi/\Sigma$ & 3  & 5  & 3 & 1.67 & 1.50 & 0.17 & 3 / 3   & 1 / 1      & --- & --- \\
5  & $S^1$       & 5   & 7  & 3 & 2.33 & 2.14 & 0.19 & 3 / 3   & 1 / 1      & --- & --- \\
6  & PropTrunc   & 8   & 8  & 3 & 2.67 & 2.56 & 0.11 & 4 / 3   & 3 / 3      & $S^2$ (3.33) & $S^1$ (2.33) \\
7  & $S^2$       & 13  & 10 & 3 & 3.33 & 3.00 & 0.33 & 4 / 3   & 2 / 2      & --- & $S^1$ (2.33) \\
8  & $S^3$       & 21  & 18 & 5 & 3.60 & 3.43 & 0.17 & 5 / 5   & 1 / 1      & --- & --- \\
9  & Hopf        & 34  & 17 & 4 & 4.25 & 4.01 & 0.24 & 5 / 4   & 1 / 1      & --- & --- \\
10 & Cohesion    & 55  & 19 & 4 & 4.75 & 4.46 & 0.29 & 5 / 4   & 2 / 2      & --- & Hopf (4.25) \\
11 & Connections & 89  & 27 & 5 & 5.40 & 4.91 & 0.49 & 6 / 5   & 2 / 2      & --- & $S^3$ (3.60) \\
12 & Curvature   & 144 & 35 & 6 & 5.83 & 5.47 & 0.36 & 6 / 5|6 & 3 / 3      & --- & $S^3$ (3.60) \\
13 & Metric      & 233 & 47 & 7 & 6.71 & 6.07 & 0.65 & 8 / 7   & 3 / 3      & --- & HIT (2.29) \\
14 & Hilbert     & 377 & 63 & 9 & 7.00 & 6.78 & 0.22 & 9 / 8|9 & 4 / 4      & --- & HIT (3.00) \\
15 & DCT         & 610 & 88 & 8 & 11.00 & 7.51 & 3.49 & 9 / 8   & 3 / 3      & --- & HIT (3.00) \\
\bottomrule
\end{tabular}
\end{table}

\subsection{Bootstrap: Steps 1--4}
\label{subsec:genesis-bootstrap-proof}

The bootstrap steps establish the minimum typed infrastructure from which later
homotopical and geometric structure can be sealed.
The foundational rule counts used here are the exact strict values recorded in
\cref{rem:automation-scope,tab:genesis-sequence-trace,tab:mbtt-audit}.

\begin{proposition}[Step 1: Universe]
\label{prop:genesis-step1}
At Step~1 the strict genesis trace selects the universe package.
\end{proposition}

\begin{proof}
\paragraph{What it is.}
The selected structure is the first universe $\U_0$, namely the ambient type of
small types.
In the strict trace it is the unique selected bootstrap object at step~1; see
\cref{tab:genesis-sequence-trace}.

\paragraph{Why $\kappa$ has this value.}
The charged cost is $\kappa=2$, exactly the two foundational clauses audited in
\cref{tab:mbtt-audit}: universe formation and the universe-level witness.
This is the minimal non-empty package that realizes a type universe rather than
an opaque constant, so the cost is fixed by \cref{def:effort}.

\paragraph{How it builds on the established library.}
It builds on nothing, because the prior library is empty.
Precisely for that reason it is the only admissible way to create a typed growth
regime at all: every later candidate presupposes some ambient type formation.

\paragraph{Why the novelty count has this value.}
Its novelty is the exact foundational count $\nu=1$ from
\cref{rem:automation-scope}: one new formation rule for small types.
There is no earlier eliminator to interact with, so no additional
cross-interface novelty is available at this step.

\paragraph{Why it is the unique selected winner at this step.}
The strict trace opens cap/band $2/2$, evaluates $2144$ raw candidates, and
selects Universe with $\rho=0.50$ exactly at the initial bar $0.50$; see
\cref{tab:genesis-sequence-trace}.
By the bootstrap admissibility conditions encoded in the strict search and by
\cref{ax:selection}, this is the unique least-supercritical bootstrap package.
\end{proof}

\begin{proposition}[Step 2: Unit]
\label{prop:genesis-step2}
At Step~2 the strict genesis trace selects the unit type $\mathbf{1}$.
\end{proposition}

\begin{proof}
\paragraph{What it is.}
The selected structure is the terminal one-point type $\mathbf{1}$.
It is the first concrete small type over $\U_0$, and therefore the first place
where inhabitation can later be witnessed.

\paragraph{Why $\kappa$ has this value.}
The charged cost is $\kappa=1$ because only the unit-type formation clause is
irreducibly new at this step; see \cref{tab:mbtt-audit}.
No extra eliminator clause is charged yet, since the selected package here is
the type former itself, not its witness.

\paragraph{How it builds on the established library.}
It leverages Step~1 exactly once: $\mathbf{1}$ is introduced as a small type in
$\U_0$.
This makes Step~2 the first genuinely inhabited target available for subsequent
introduction rules.

\paragraph{Why the novelty count has this value.}
The strict total is again exact and equals $\nu=1$; see
\cref{rem:automation-scope,tab:genesis-sequence-trace}.
At this stage the unit type contributes one new foundational rule, but still no
non-trivial interface interaction.

\paragraph{Why it is the unique selected winner at this step.}
The strict trace opens cap/band $1/1$ and selects Unit at
$\rho=1.00 > 0.50$, with no displayed runner-up or near miss; see
\cref{tab:genesis-sequence-trace}.
Since the winner clears the bar with the smallest available cost and bootstrap
admissibility still restricts the search to one-clause packages over the new
universe, \cref{ax:selection} forces this choice.
\end{proof}

\begin{proposition}[Step 3: Witness]
\label{prop:genesis-step3}
At Step~3 the strict genesis trace selects the unit witness
$\star:\mathbf{1}$.
\end{proposition}

\begin{proof}
\paragraph{What it is.}
The selected structure is the canonical inhabitant of the unit type.
Semantically it is the first realized term, not merely the first realized type.

\paragraph{Why $\kappa$ has this value.}
The charged cost is $\kappa=1$, because the selected specification adds exactly
one irreducible introduction clause, namely the witness itself; see
\cref{tab:mbtt-audit}.
The associated elimination behavior is not charged separately at this step
because it is accounted for via adjoint completion rather than as an additional
primitive clause.

\paragraph{How it builds on the established library.}
It uses the already selected unit type as its codomain and thus turns the
bootstrap library from a pure type-formation fragment into an inhabited theory.
This is exactly the bridge from Step~2's abstract type to a concrete term.

\paragraph{Why the novelty count has this value.}
By \cref{lem:adjoint-completion,rem:witness-pisigma-attribution}, the witness
contributes one introduction rule and one elimination action, giving the exact
strict total $\nu=2$ recorded in \cref{tab:genesis-sequence-trace}.
This is the canonical early example where eliminative payload exceeds the raw
charged clause count.

\paragraph{Why it is the unique selected winner at this step.}
The strict trace opens cap/band $1/1$ and selects Witness with
$\rho=2.00$ against bar $1.33$; the displayed top near miss is only a
foundational candidate at $\rho=1.00$; see
\cref{tab:genesis-sequence-trace}.
Hence \cref{ax:selection} picks the witness package as the unique bar-clearer
with minimal overshoot in the foundational one-clause regime.
\end{proof}

\begin{proposition}[Step 4: Dependent function and sum types]
\label{prop:genesis-step4}
At Step~4 the strict genesis trace selects the dependent function/sum package
($\Pi/\Sigma$ types).
\end{proposition}

\begin{proof}
\paragraph{What it is.}
The selected structure is the minimal former package supplying dependent
functions, dependent sums, and the core application/pairing interface.
Although the engine names the winner ``Pi'' in the trace, the mathematical
structure is the joint $\Pi/\Sigma$ former package described in
\cref{tab:mbtt-audit}.

\paragraph{Why $\kappa$ has this value.}
The charged cost is $\kappa=3$, corresponding to the selected specification
listed in \cref{tab:mbtt-audit}: $\lambda$-introduction, pair formation, and
application.
This is the minimal cost at which the library acquires bona fide dependent
former structure.

\paragraph{How it builds on the established library.}
It builds on the previous three bootstrap entries by turning the bare inhabited
universe into a theory with dependent abstraction and pairing.
That infrastructure is the prerequisite for all later HIT, map, and API-shell
steps.

\paragraph{Why the novelty count has this value.}
By \cref{rem:witness-pisigma-attribution}, the exact foundational total is
$\nu=5$: two introduction rules (abstraction and pairing) together with three
elimination rules (application, first projection, second projection).
This is the strict total reported in \cref{tab:genesis-sequence-trace}.

\paragraph{Why it is the unique selected winner at this step.}
The strict trace opens cap/band $3/3$ and finds a singleton viable set, with
$\rho=1.67$ against bar $1.50$; see \cref{tab:genesis-sequence-trace}.
Since Step~4 is the first stage where the library needs a genuine former layer,
and since the selected package is already constructively prime in the sense of
\cref{ax:selection,def:minimal-api}, it is uniquely selected.
\end{proof}

\subsection{Geometric Ascent: Steps 5--9}
\label{subsec:genesis-geometric-proof}

The next five steps move from bootstrap syntax to genuine higher structure.
Here the key theorem-level inputs are the two-layer coherence results
\cref{lem:trace,thm:scaling}, the HIT/map exactness status from
\cref{rem:automation-scope}, and the Agda abstraction-barrier evidence reported
in \cref{sec:verification}.

\begin{proposition}[Step 5: The circle]
\label{prop:genesis-step5}
At Step~5 the strict genesis trace selects the circle $S^1$.
\end{proposition}

\begin{proof}
\paragraph{What it is.}
The selected structure is the first higher inductive type, the circle $S^1$,
with a point and a non-trivial loop.
This is the library's first genuinely homotopical object.

\paragraph{Why $\kappa$ has this value.}
The charged cost is $\kappa=3$, because the selected conceptual shell is the
minimal three-clause HIT signature from \cref{tab:mbtt-audit}: formation,
base point, and loop.
The compiled telescope is larger in the strict trace, but by the proof contract
that extra compiled payload is not extra charged cost.

\paragraph{How it builds on the established library.}
It is the first higher object that can be expressed only after the
$\Pi/\Sigma$ former layer has been installed.
Thus it turns the bootstrap library into a setting where non-trivial path data
can be generated and reused.

\paragraph{Why the novelty count has this value.}
By \cref{ex:foundation-steps}, the exact total is $\nu=7$:
five introduction-style schemas together with two computation/path rules coming
from the loop constructor.
This is the strict HIT reconstruction recorded in
\cref{rem:automation-scope,tab:genesis-sequence-trace}.

\paragraph{Why it is the unique selected winner at this step.}
The strict trace opens cap/band $3/3$ and the viable set is singleton, with
$\rho=2.33$ against bar $2.14$; see \cref{tab:genesis-sequence-trace}.
By \cref{lem:genesis-bar-lemma,ax:selection}, the circle is therefore the
unique least-supercritical entry into the geometric regime.
\end{proof}

\begin{proposition}[Step 6: Propositional truncation]
\label{prop:genesis-step6}
At Step~6 the strict genesis trace selects propositional truncation.
\end{proposition}

\begin{proof}
\paragraph{What it is.}
The selected structure is the propositional truncation operator
$\|{-}\|_0$, which collapses higher structure to a mere proposition while still
remembering existence.
This is the first operation that systematically controls homotopical complexity
rather than only creating it.

\paragraph{Why $\kappa$ has this value.}
The charged cost is again $\kappa=3$, with the conceptual shell audited in
\cref{tab:mbtt-audit} as truncation formation, truncation introduction, and the
squash witness.
As with $S^1$, the strict compiled telescope is larger than the charged
conceptual shell, but the selector charges only the conceptual package.

\paragraph{How it builds on the established library.}
It builds directly on the new geometric branch opened by $S^1$.
In the library-growth dynamics this is the first point where PEN can both
generate higher data and quotient it to mere existence, which is why the
subsequent lift to $S^2$ becomes available.

\paragraph{Why the novelty count has this value.}
The strict total is the exact HIT reconstruction $\nu=8$ recorded in
\cref{rem:automation-scope}.
Operationally this counts the new truncation former, its introduction and
elimination payload, and the squash-induced computation structure physically
present in the strict HIT compilation.

\paragraph{Why it is the unique selected winner at this step.}
This is the tightest selection point in the whole trajectory.
The strict trace opens cap/band $4/3$ and finds exactly two viable candidates:
Trunc at $\rho=8/3$ and $S^2$ at $\rho=10/3$; see
\cref{tab:genesis-sequence-trace}.
Since both have the same charged $\kappa=3$, \cref{ax:selection} reduces to
minimal overshoot, and Trunc wins because its margin is only $0.11$ whereas
$S^2$ overshoots by about $0.77$.
The worked arithmetic in \cref{app:worked-arithmetic} shows that even one extra
charged clause would kill this step, so the winner is both unique and
structurally critical.
\end{proof}

\begin{proposition}[Step 7: The 2-sphere]
\label{prop:genesis-step7}
At Step~7 the strict genesis trace selects the sphere $S^2$.
\end{proposition}

\begin{proof}
\paragraph{What it is.}
The selected structure is the first genuinely two-dimensional sphere, realized
as a minimal higher inductive package.
It is the next geometric lift after the circle/truncation bridge has been
installed.

\paragraph{Why $\kappa$ has this value.}
The charged cost remains $\kappa=3$, because the selected conceptual shell is
the three-clause $S^2$ signature from \cref{tab:mbtt-audit}: formation, base
point, and 2-path surface witness.
Again the strict selector charges the conceptual shell, not every derived
compiled clause.

\paragraph{How it builds on the established library.}
By \cref{lem:trace,thm:scaling}, Step~7 inherits the two-layer interface debt
created by Steps~5 and~6.
The geometric meaning is exactly that $S^2$ is the first post-trunc higher lift:
the library now has both non-trivial loop data and a way to compress it.

\paragraph{Why the novelty count has this value.}
The strict total is the exact HIT reconstruction $\nu=10$ listed in
\cref{rem:automation-scope,tab:genesis-sequence-trace}.
It is larger than the circle's total because a 2-dimensional HIT carries a
richer path/computation payload in the strict evaluator.

\paragraph{Why it is the unique selected winner at this step.}
The strict trace opens cap/band $4/3$ and the viable set is singleton:
$S^2$ clears the bar at $\rho=3.33$, while the displayed near miss is only the
replayed $S^1$ package at $\rho=2.33$; see
\cref{tab:genesis-sequence-trace}.
Thus \cref{ax:selection} again yields a unique winner.
\end{proof}

\begin{proposition}[Step 8: The H-space on $S^3$]
\label{prop:genesis-step8}
At Step~8 the strict genesis trace selects the H-space package on $S^3$.
\end{proposition}

\begin{proof}
\paragraph{What it is.}
The selected structure is not merely a bare 3-sphere, but the minimal
multiplicative H-space package on $S^3$.
This is the first point where PEN prefers a slightly richer geometric shell to a
bare sphere because the additional interaction is worth the extra cost.

\paragraph{Why $\kappa$ has this value.}
The charged cost is $\kappa=5$, exactly the H-space specification selected in
\cref{tab:mbtt-audit}.
The paper's comparison in \cref{rem:kolmogorov-ambiguity} explains why this is
the relevant charged shell: bare 3-sphere encodings exist at lower $\kappa$, but
they overshoot the bar too much and are therefore not selected.

\paragraph{How it builds on the established library.}
By \cref{lem:trace}, Step~8 resolves the Fibonacci debt
$21=13+8$ inherited from Steps~7 and~6.
Geometrically, it uses the just-installed $S^2$ layer together with the
truncation bridge to support the first multiplicative higher sphere.

\paragraph{Why the novelty count has this value.}
The strict total is the exact HIT/H-space reconstruction $\nu=18$ from
\cref{rem:automation-scope}.
More specifically, \cref{rem:kolmogorov-ambiguity} shows that the bare 3-sphere
already contributes $15$ rules, and the H-space laws add $3$ further
elimination-style rules, yielding $18$ in total.

\paragraph{Why it is the unique selected winner at this step.}
The verified strict trace already has a singleton viable set at cap/band $5/5$;
see \cref{tab:genesis-sequence-trace}.
The paper's stronger specification comparison in
\cref{rem:kolmogorov-ambiguity} then explains the uniqueness semantically:
among the natural $S^3$ presentations that clear the bar, minimal overshoot
selects the H-space package rather than the bare suspension or the bare native
HIT.
\end{proof}

\begin{proposition}[Step 9: The Hopf fibration]
\label{prop:genesis-step9}
At Step~9 the strict genesis trace selects the Hopf fibration.
\end{proposition}

\begin{proof}
\paragraph{What it is.}
The selected structure is the map package
$h:S^3\to S^2$ classifying a principal $S^1$-bundle over $S^2$.
This is the first step where the genesis sequence crosses from higher types to a
non-trivial higher map.

\paragraph{Why $\kappa$ has this value.}
The charged cost is $\kappa=4$, exactly the selected Hopf specification audited
in \cref{tab:mbtt-audit}: map, fiber, total-space witness, and classifying map.
Unlike the earlier HIT steps, the conceptual shell and the charged shell agree
here.

\paragraph{How it builds on the established library.}
It leverages the full preceding homotopy package:
the domain $S^3$, the codomain $S^2$, and the fiber $S^1$ are all earlier
library entries.
By \cref{lem:trace} and the Agda evidence summarized in
\cref{sec:verification}, the inherited obligation count is exactly
$34=21+13$, with a domain/codomain split reflecting the $S^3/S^2$ interface.

\paragraph{Why the novelty count has this value.}
For Step~9 the paper gives a sharp qualitative description:
\cref{sec:discussion} records that $\nu_G=\nu_H=0$ and $\nu=\nu_C$.
Thus the exact strict total $\nu=17$ from
\cref{rem:automation-scope,tab:genesis-sequence-trace} is entirely interactional
map novelty rather than new type-former novelty.

\paragraph{Why it is the unique selected winner at this step.}
The strict trace opens cap/band $5/4$ and the viable set is singleton, so Hopf
is already the unique strict bar-clearer at $\rho=4.25$ against bar $4.01$; see
\cref{tab:genesis-sequence-trace}.
This is the point where the codebase's structural regime shifts from higher
types to higher maps, but in the canonical strict lane no additional rescue
heuristic is needed: the viable set has already collapsed to a single winner.
\end{proof}

\subsection{Framework Abstraction: Steps 10--14}
\label{subsec:genesis-framework-proof}

From Step~10 onward the sequence grows by library specifications rather than by
new foundational judgment forms; see \cref{rem:library-api}.
The semantic force of each step comes from constructive irreducibility and
structural unity, not from opaque axiom packing.
For Step~10 the rule audit is exact; for Steps~11--14 it is expository but
constrained to match the strict totals by \cref{rem:automation-scope}.

\begin{proposition}[Step 10: Cohesion]
\label{prop:genesis-step10}
At Step~10 the strict genesis trace selects the cohesive modality package.
\end{proposition}

\begin{proof}
\paragraph{What it is.}
The selected structure is the adjoint quadruple of modalities
$(\flat,\sharp,\Disc,\Pi_{\infty})$ described in
\cref{app:cohesion-spec}.
This is the first abstraction step: PEN stops adding one more homotopy object
and instead adds a modal interface for the whole geometric library.

\paragraph{Why $\kappa$ has this value.}
The charged cost is $\kappa=4$, exactly the four modality-formation clauses
listed in \cref{app:cohesion-spec,tab:mbtt-audit}.
By \cref{def:derivability,def:minimal-api}, these are irreducible API clauses:
the adjoint string cannot be obtained definitionally from the prior HoTT
library.

\paragraph{How it builds on the established library.}
By \cref{rem:library-api}, a late API shell is valuable only insofar as it
interacts with earlier structure.
That interaction is explicit here:
\cref{app:cohesion-spec} records cross-interactions with $S^1$, $S^2$, $S^3$,
and Hopf, so Cohesion is a global abstraction of the entire geometric ascent.

\paragraph{Why the novelty count has this value.}
For Step~10 the rule audit is exact:
\cref{app:cohesion-spec} gives $\nu=19=2+17$, namely two introduction-style unit
maps together with seventeen elimination/interaction rules.
This exact total is the one selected in the strict trace of
\cref{tab:genesis-sequence-trace}.

\paragraph{Why it is the unique selected winner at this step.}
The strict trace opens cap/band $5/4$ and again finds a singleton viable set:
Cohesion clears the bar at $\rho=4.75$, while the displayed near miss is the old
Hopf package at $\rho=4.25$; see \cref{tab:genesis-sequence-trace}.
Hence \cref{ax:selection} selects the cohesive shell uniquely, and the
non-survival of Hopf shows that PEN is now in a genuinely new abstraction
regime.
\end{proof}

\begin{proposition}[Step 11: Connections]
\label{prop:genesis-step11}
At Step~11 the strict genesis trace selects the connection package.
\end{proposition}

\begin{proof}
\paragraph{What it is.}
The selected structure is the differential-geometric package centered on a
connection $\nabla$ together with transport, covariant lift, horizontal lift,
and Leibniz interaction; see \cref{app:connections-spec}.
This is the first differential shell built over the cohesive interface.

\paragraph{Why $\kappa$ has this value.}
The charged cost is $\kappa=5$, exactly the five API clauses listed in
\cref{app:connections-spec,tab:mbtt-audit}.
By \cref{def:derivability}, a connection is not derivable from cohesion alone,
so the package must be charged explicitly.

\paragraph{How it builds on the established library.}
It leverages Cohesion as its immediate carrier and thereby translates the modal
library into differential transport data.
This is the first late step whose entire meaning is cross-interface leverage:
without the cohesive layer from Step~10 there is no differential package to
select.

\paragraph{Why the novelty count has this value.}
The expository rule audit in \cref{app:connections-spec} reconstructs the strict
total as $\nu=27=2+25$, where the new interaction comes from cohesive
cross-rules, function-space rules, and library interactions with the previously
selected HITs and Hopf.
By \cref{rem:automation-scope}, this reconstruction is subordinate to, but
consistent with, the strict total in \cref{tab:genesis-sequence-trace}.

\paragraph{Why it is the unique selected winner at this step.}
The strict trace opens cap/band $6/5$ and the viable set is singleton:
Connections clears the bar at $\rho=5.40$, while the displayed near miss is only
the old $S^3$ package at $\rho=3.60$; see
\cref{tab:genesis-sequence-trace}.
Thus the strict lane does not need any external steering here: the first
differential shell is already the unique bar-clearing continuation.
\end{proof}

\begin{proposition}[Step 12: Curvature]
\label{prop:genesis-step12}
At Step~12 the strict genesis trace selects the curvature package.
\end{proposition}

\begin{proof}
\paragraph{What it is.}
The selected structure is the curvature layer
$R=d\nabla+\nabla\wedge\nabla$ together with Bianchi, holonomy, Chern--Weil, and
characteristic-class structure; see \cref{app:curvature-spec}.
It is the first shell whose meaning is ``second-order differential structure''
rather than just differential transport.

\paragraph{Why $\kappa$ has this value.}
The charged cost is $\kappa=6$, exactly the six irreducible clauses audited in
\cref{app:curvature-spec,tab:mbtt-audit}.
These clauses are constructively irreducible because the elimination behavior
they expose is not definitionally recoverable from a bare connection package;
see \cref{def:derivability,app:curvature-spec}.

\paragraph{How it builds on the established library.}
It leverages Connections as its immediate carrier and Cohesion as its ambient
modal base.
This is the canonical next abstraction after Step~11: once the library can
differentiate, the next non-derivable interaction is precisely the failure of
flatness.

\paragraph{Why the novelty count has this value.}
The expository audit in \cref{app:curvature-spec} reconstructs the strict total
as $\nu=35=2+33$, with the extra weight coming from interactions among
curvature, connection, and cohesive structure.
By \cref{rem:automation-scope}, that reconstruction is explanatory rather than a
second oracle, but it agrees with the canonical strict total in
\cref{tab:genesis-sequence-trace}.

\paragraph{Why it is the unique selected winner at this step.}
The strict trace opens cap/bands $6/(5|6)$ and nevertheless yields a singleton
viable set: Curvature clears the bar at $\rho=5.83$, while the displayed near
misses are far below the bar; see \cref{tab:genesis-sequence-trace}.
Thus the switch from differential to curvature structure is not a naming
convention but a mechanically unique strict continuation.
\end{proof}

\begin{proposition}[Step 13: Metric plus frame bundle]
\label{prop:genesis-step13}
At Step~13 the strict genesis trace selects the metric package rather than a
bare tensor or scalar-heavy detour.
\end{proposition}

\begin{proof}
\paragraph{What it is.}
The selected structure is the metric-plus-frame-bundle package described in
\cref{app:metric-spec,tab:step13}: a symmetric bilinear form together with
Levi-Civita connection, geodesic, volume, Hodge star, Laplacian, and Ricci
structure.
This is the first step where PEN selects an explicitly geometric API bundle
rather than a single operator.

\paragraph{Why $\kappa$ has this value.}
The charged cost is $\kappa=7$, exactly the seven clauses audited in
\cref{app:metric-spec,tab:step13}.
By \cref{def:minimal-api}, this package is charged as a single structural bundle:
the bare metric tensor is too weak, whereas the coupled bundle is the smallest
constructively prime package that can survive the Step~13 bar.

\paragraph{How it builds on the established library.}
It leverages Curvature and Connections directly and uses the cohesive/differential
layers indirectly through those references.
The key structural point, emphasized in \cref{app:metric-spec}, is that the
metric package is selected because it unlocks the whole chain
$g\to\nabla_g\to R\to\mathrm{Ric}$ inside the already accumulated differential
library.

\paragraph{Why the novelty count has this value.}
The expository reconstruction in \cref{app:metric-spec} gives
$\nu=47=4+43$, where the dominant term is the bounded coupling with the
curvature and connection layers.
The same section explains why a bare metric tensor would have
$\rho\le 3.0$, far below the Step~13 bar, so the strict total in
\cref{tab:genesis-sequence-trace} is genuinely bundle-driven rather than
tensor-driven.

\paragraph{Why it is the unique selected winner at this step.}
The strict trace opens cap/band $8/7$ and the viable set is singleton:
Metric clears the bar at $\rho=6.71$, whereas the displayed near miss is only an
unnamed HIT candidate at $\rho=2.29$; see
\cref{tab:genesis-sequence-trace}.
Hence the selected metric bundle is uniquely forced by the strict objective, and
\cref{def:minimal-api} explains why smaller fragments are excluded.
\end{proof}

\begin{proposition}[Step 14: Hilbert functional]
\label{prop:genesis-step14}
At Step~14 the strict genesis trace selects the Hilbert-functional package.
\end{proposition}

\begin{proof}
\paragraph{What it is.}
The selected structure is the operator-analytic package described in
\cref{app:hilbert-spec,tab:step14}: inner product, completeness, orthogonal and
spectral decomposition, $C^*$-algebra structure, cross-links to Metric,
Curvature, and Connections, and a functional derivative.
This is the analytic culmination of the framework-abstraction phase.

\paragraph{Why $\kappa$ has this value.}
The charged cost is $\kappa=9$, exactly the nine clauses audited in
\cref{app:hilbert-spec,tab:step14}.
By \cref{def:derivability}, none of this operator-analytic structure is
definitionally free from the preceding geometric library.

\paragraph{How it builds on the established library.}
It leverages the three immediately relevant layers: Metric, Curvature, and
Connections.
This is a textbook PEN abstraction step: the new package is not selected for
opaque existence but because it compresses and re-exposes the interaction of the
whole geometric stack in analytic form.

\paragraph{Why the novelty count has this value.}
The expository reconstruction in \cref{app:hilbert-spec} gives
$\nu=63=5+58$, dominated by the coupling to Metric, Curvature, and Connections.
The same section also shows why the margin is small: even two extra charged
clauses would drop the package below the Step~14 bar.

\paragraph{Why it is the unique selected winner at this step.}
The strict trace opens cap/bands $9/(8|9)$ and still yields a singleton viable
set: Hilbert clears the bar at $\rho=7.00$, while the displayed near miss is
only an unnamed HIT candidate at $\rho=3.00$; see
\cref{tab:genesis-sequence-trace}.
Thus the operator-analytic layer is the unique strict continuation of the metric
bundle, with no competing bar-clearer in the verified run.
\end{proof}

\subsection{Terminal Synthesis: Step 15}
\label{subsec:genesis-terminal-proof}

Step~15 is special.
The mathematical interpretation as DCT is theorem-level and semantic, but the
only canonical quantitative claim is the strict total $\nu=88$;
see \cref{sec:decomposition,rem:automation-scope}.

\begin{proposition}[Step 15: Temporal-cohesive synthesis]
\label{prop:genesis-step15}
At Step~15 the strict genesis trace selects the temporal-cohesive shell whose
semantic completion is read as DCT.
\end{proposition}

\begin{proof}
\paragraph{What it is.}
The selected structure is the temporal-cohesive extension audited in
\cref{tab:step15}: next modality, eventually modality, the map
$\bigcirc\to\Diamond$, and the explicit exchange laws between temporal and
cohesive structure.
Its semantic completion is the DCT reading established in
\cref{sec:decomposition}.

\paragraph{Why $\kappa$ has this value.}
The charged cost is $\kappa=8$, exactly the eight temporal-cohesive clauses
audited in \cref{tab:step15}.
As stressed in \cref{tab:step15,sec:decomposition}, the infinitesimal reading is
\emph{not} charged as an extra primitive clause; it is derived semantically from
the temporal-cohesive shell already selected.

\paragraph{How it builds on the established library.}
It builds directly on Cohesion and sits atop the already selected Hilbert layer.
This is the terminal synthesis step because it fuses the earlier spatial logic
of Step~10 with temporal operators in a way that internalizes the library's own
evolution; see \cref{thm:end-of-history,sec:decomposition}.

\paragraph{Why the novelty count has this value.}
Only the strict total $\nu=88$ is canonical here; no finer split is asserted in
this section.
By \cref{sec:verification,sec:adjoint-completion}, the strict engine computes
that total from the conceptual temporal shell together with the explicit
path-style witness payload preserved by compilation.
This is exactly the number reported in \cref{tab:genesis-sequence-trace}.

\paragraph{Why it is the unique selected winner at this step.}
The strict trace opens cap/band $9/8$ and the viable set is singleton:
DCT clears the bar at $\rho=11.00$ against bar $7.51$, while the displayed near
miss is only an unnamed HIT candidate at $\rho=3.00$; see
\cref{tab:genesis-sequence-trace}.
This is also the largest absolute overshoot in the entire sequence.
Thus the temporal-cohesive shell is not merely admissible but uniquely dominant
within the strict lane at the terminal step.
\end{proof}

\subsection{Conclusion of the Step-by-Step Proof}
\label{subsec:genesis-proof-conclusion}

\begin{theorem}[The strict genesis sequence]
\label{thm:strict-genesis-sequence}
In the current strict unguided PEN lane, the unique step-by-step sequence of
selected structures through Step~15 is
\[
  \begin{aligned}
    &\U_0,\ \mathbf{1},\ \star,\ \Pi/\Sigma,\ S^1,\ \|{-}\|_0,\ S^2,\ S^3,\\
    &\mathrm{Hopf},\ \mathrm{Cohesion},\ \mathrm{Connections},\\
    &\mathrm{Curvature},\ \mathrm{Metric},\ \mathrm{Hilbert},\\
    &\mathrm{DCT}.
  \end{aligned}
\]
For each $1\le n\le 15$, the selected structure has the corresponding values of
$\Delta_n$, $\nu_n$, $\kappa_n$, $\rho_n$, and $\mathrm{Bar}_n$ listed in
\cref{tab:genesis-sequence-trace}.
\end{theorem}

\begin{proof}
Combine the global bar formula of \cref{lem:genesis-bar-lemma} with the fifteen
step propositions
\cref{prop:genesis-step1,prop:genesis-step2,prop:genesis-step3,prop:genesis-step4,prop:genesis-step5,prop:genesis-step6,prop:genesis-step7,prop:genesis-step8,prop:genesis-step9,prop:genesis-step10,prop:genesis-step11,prop:genesis-step12,prop:genesis-step13,prop:genesis-step14,prop:genesis-step15}.
Each step proposition identifies the selected structure, justifies its charged
cost and novelty total at the appropriate proof-strength level, and verifies
that it is the unique strict-lane winner at its step.
\end{proof}

\begin{corollary}[No Step~16 continuation in the current strict lane]
\label{cor:genesis-no-step16}
No Step~16 candidate clears the selection bar in the current strict unguided
lane.
\end{corollary}

\begin{proof}
By \cref{thm:strict-genesis-sequence}, Step~15 is the selected temporal-cohesive
shell with the canonical strict total $\nu=88$.
By the halting theorem \cref{thm:end-of-history}, every internal Step~16
continuation has zero marginal external novelty and every disconnected
continuation fails the positive Step~16 bar.
Hence no Step~16 candidate clears the bar.
\end{proof}

\let\paragraph\genesissequenceorigparagraph

\ifgenesissequencestandalone
\end{document}
\fi
